\chapter{Phase 3: the interaction term}\label{ch:cII3}
\chapterbrief
{Part~I, Chapter~\ref{ch:cI3} (A19--A24): the class $\Bclass$, the closed form, the cost, the instability, Dyson--Phillips and the no-go for infinite adjoint orbits.}
{to prove the supporting lemmas of the class $\Bclass$, the results on factors and second-order evolutions inside $\Bclass$, the solvable and non-solvable Lie cases, the interaction created by a boundary, the exact frontier between continuous and discontinuous interaction, the balance decomposition, and the residual-based no-go results.}
{Dim.~B43--B60.}

%=====================================================================
\section{The class $\Bclass$}

\begin{dimcard}{B43}{P1-L1}{Properties of $\Bclass$}{\proved}{def:B}
(a) $B_\rho$, $\Bclass_\rho$ are Banach spaces and $\|\partial_x\|_{B_\rho\to B_\rho}\le\rho$. (b) $\Bclass_\rho\hookrightarrow H^m$ for every $m$. (c) Every exponential-polynomial lies in $\Bclass$. (d) For $u\in\Bclass_\rho$, $t\mapsto u(t,\cdot)$ extends to $\bar I$ and lies in $C(\bar I;B_{\rho'})$ for $\rho'>\rho$.
\end{dimcard}
\begin{proof}
\step{1}{(a)}
A Cauchy sequence has all derivatives Cauchy in $L^2$; the limits satisfy $g_{j+1}=g_j'$, so $g_0\in H^\infty$ with $g_0^{(j)}=g_j$, uniformly in $j$. Shifting the index gives $\|f'\|_{B_\rho}\le\rho\|f\|_{B_\rho}$.
\step{2}{(b), (c)}
$\|u\|_{H^m}^2\le\binom{m+2}{2}\max(1,\rho)^{2m}\|u\|_{\Bclass_\rho}^2$. For $x^le^{\zeta x}$ the $j$-th derivative is $e^{\zeta x}$ times a polynomial with coefficients $O(j^l|\zeta|^j)$, and $j^l(|\zeta|/\rho)^j$ is bounded for $\rho>|\zeta|$.
\step{3}{(d)}
$H^1(I;L^2(J))\hookrightarrow C(\bar I;L^2(J))$ applied to $\partial_x^ju$, with the geometric decay $(\rho/\rho')^j$.
\whyworks{$\Bclass$ is the space where derivatives grow at most geometrically, so exponentials of polynomials in $\partial_x$ converge.}
\end{proof}
\uses{none.}

\begin{dimcard}{B44}{P1-L2}{Flows of terms in $\Bclass$}{\proved}{}
On $B_\rho$: $\|A(\partial_x)\|\le|A|(\rho)$; $e^{sA(\partial_x)}$ converges in norm with $\|e^{sA}\|\le e^{|s||A|(\rho)}$; the family is entire in $s$, a group, independent of $\rho$; exponentials of different polynomials commute and $e^{s(A+A')}=e^{sA}e^{sA'}$; and $e^{sA(\partial_x)}(x^le^{\zeta x})=\partial_\zeta^l\big(e^{sA(\zeta)}e^{\zeta x}\big)$.
\end{dimcard}
\begin{proof}
\step{1}{Norm and commutation}
$\|A(\partial)\|\le\sum|\alpha_l|\|\partial\|^l$. Polynomials in one operator commute, and in a Banach algebra $e^{X+Y}=e^Xe^Y$ for commuting $X,Y$ (Cauchy product).
\step{2}{Action on modes}
$\zeta\mapsto e^{\zeta x}\in B_\rho$ is holomorphic for $|\zeta|<\rho$ and $A(\partial)^ke^{\zeta x}=A(\zeta)^ke^{\zeta x}$; a bounded operator commutes with $\partial_\zeta^l$.
\whyworks{on $\Bclass$ the derivative is a bounded operator, so its polynomials have exponentials.}
\end{proof}
\uses{B43.}

\begin{dimcard}{B45}{P1-L3}{The Cauchy problem in the class}{\proved}{}
\phantomsection\label{dim:P1-L4}\phantomsection\label{dim:P1-L5}%
(P1-L3) A continuous curve $v$ in $B_\rho$ with $(\partial_t-A(\partial_x))v=0$ is $v(t)=e^{(t-t_0)A}v(t_0)$; conversely these are solutions in $\Bclass$. (P1-L4) For $r\in\C[\xi]$ the Volterra primitive $w=\int_{x_0}^xG(x-y)F(t,y)\dd y$ ($G$ the fundamental solution of $r(\partial)$) satisfies $r(\partial_x)w=F$ and maps $\Bclass_\rho$ into $\Bclass_R$. (P1-L5) Formal lemmas in $\mathcal A[[t]]$: $e^{gX}Ye^{-gX}=\sum_kg^k\mathrm{ad}_X^kY/k!$, and for $U=\prod_ie^{g_iX_i}$, $U'U^{-1}=\sum_ig_i'\,\mathrm{Ad}(P_{i-1})X_i$.
\end{dimcard}
\begin{proof}
\step{1}{(L3)}
Testing the equation against $\varphi\otimes\psi$ shows that the distributional time derivative of $v$, as a $B_\rho$-valued distribution, is $A(\partial)v$; a Banach-valued distribution with zero derivative is constant, so $\frac{d}{dt}\big(e^{-(t-t_0)A}v\big)=0$.
\step{2}{(L4)}
Variation of constants gives $r(\partial_x)w=F$. For $j\ge\delta=\deg r$, $\partial_x^jw$ is expressed through $\partial_x^{j-\delta}F$ and lower derivatives of $w$; induction gives $\sup_i\rho^{-i}\|\partial_t^i\partial_x^jw\|\le KR^j$.
\step{3}{(L5)}
Both sides of each identity solve the same formal ODE with the same initial value.
\whyworks{uniqueness in $\Bclass$ is what excludes Tychonoff and boundary-fed solutions.}
\end{proof}
\uses{B43, B44.}

\begin{dimcard}{B46}{P1-P1}{Single flows are not determined without a class}{\proved{} / \cf{} (Tychonoff)}{prop:nonuniq}
(a) For $a_n=\xi$ the solutions with a given datum on $J$ are $F(x+t)$ for any extension $F$: infinitely many. (b) For $a_n=\xi^2$ there are nonzero heat solutions with zero datum; those with datum $u_0$ form an infinite-dimensional affine space. (c) So without a class $\Eint$ is not a function of $u$; in $\Bclass$ it is unique.
\end{dimcard}
\begin{proof}
\step{1}{(a)}
In coordinates $(t,y=x+t)$ the equation is $\partial_tw=0$ on a parallelogram, so $w=F(y)$, fixed on $J$ and free elsewhere.
\step{2}{(b)}
Tychonoff's $v(t,x)=\sum_kg^{(k)}(t)\,x^{2k}/(2k)!$ with $g(t)=e^{-t^{-2}}$ for $t>0$, $0$ otherwise, is smooth and vanishes for $t\le0$ \cf. Its time shifts $v(t-\sigma,x)$ are independent: on the first interval only the earliest survives, and $v(s,0)=g(s)\ne0$.
\step{3}{(c)}
B45(L3).
\whyworks{on a bounded domain both the transport and the heat flow need information that the datum does not contain.}
\end{proof}
\uses{B45.}

%=====================================================================
\section{Factors and second-order evolutions in $\Bclass$}

\begin{dimcard}{B47}{P1-P2}{Factors in $\Bclass$}{\proved}{}
\phantomsection\label{dim:P1-P3}%
On a bounded rectangle, $u\in\Bclass\cap\ker P$: (a) if the factor curves are disjoint, the components $E_j(\partial)u$ lie in $\Bclass$; (b) if one decomposition lies in $\Bclass$, all do; (c) for $P=\prod_j(\tau-a_j(\xi))$, every decomposition lies in $\Bclass$. \textbf{(P1-P3)} The interaction per term inside each factor is then defined modulo a finite-dimensional space of exponential-polynomials, of dimension at most $\delta(P)$.
\end{dimcard}
\begin{proof}
\step{1}{(a), (b)}
$\Bclass$ is closed under derivatives; the ambiguity consists of exponential-polynomials, which lie in $\Bclass$.
\step{2}{(c)}
By B27, $\Delta(\partial_x)u_j=F_j\in\Bclass$. B45(L4) gives $w\in\Bclass$ with $\Delta(\partial_x)w=F_j$; then $h=u_j-w\in\ker\Delta(\partial_x)=\mathcal D'(I)\otimes V_\Delta$, and $(\partial_t-a_j)h\in\Bclass$ turns the coefficients into solutions of a linear ODE system in $t$ with $\Bclass$ data, hence of geometric growth.
\step{3}{(P3)}
Changing decomposition adds $(w_j)\in\mathcal N$, which changes each interaction term linearly by an exponential-polynomial.
\whyworks{the Vandermonde formula expresses each component through derivatives of $u$, up to an ODE in $x$ that preserves geometric growth.}
\end{proof}
\uses{B27, B45.}

\begin{dimcard}{B48}{P1-P4}{Second-order factors: propagator and Gauss form}{\proved}{}
\phantomsection\label{dim:II-T7}%
For $Q=\tau^2+b(\xi)\tau+c(\xi)$ and $u\in\Bclass_\rho\cap\ker Q$: (a) $U=(u,\partial_tu)$ satisfies $U(t)=e^{(t-t_0)C(\partial_x)}U(t_0)$, with $e^{tC}$ given by the canonical fundamental system $\varphi_1,\varphi_2$ of $\varphi''+b(\zeta)\varphi'+c(\zeta)\varphi=0$; (b) for $|t|<T^*(\rho')$, with $T^*\ge\log2/K(\rho')$, the Gauss factorisation
\[
e^{tC}=\begin{pmatrix}1&0\\\varphi_1'/\varphi_1&1\end{pmatrix}\begin{pmatrix}\varphi_1&0\\0&e^{-tb}/\varphi_1\end{pmatrix}\begin{pmatrix}1&\varphi_2/\varphi_1\\0&1\end{pmatrix}
\]
holds. \textbf{(II-T7)} Mode by mode, the Wei--Norman (Gauss) form is global iff $\varphi_1(t;\zeta)$ has no zero.
\end{dimcard}
\begin{proof}
\step{1}{(a)}
B45(L3) for systems; the entries of $e^{tC(\zeta)}$ solve the scalar ODE in $t$ with the canonical initial values.
\step{2}{(b)}
Multiply out the three matrices and use the Wronskian $\varphi_1\varphi_2'-\varphi_1'\varphi_2=e^{-tb}$. An entire function without zeros on $|\zeta|\le\rho'$ is invertible in the closed algebra generated by $\partial_x$; $\varphi_1(t;\zeta)=1+O(t^2)$ has no zero for $|t|\le\log2/K$.
\whyworks{a $2\times2$ matrix factors into lower, diagonal and upper parts exactly where its top-left entry does not vanish.}
\end{proof}
\uses{B45.}

%=====================================================================
\section{Non-commuting terms}

\begin{dimcard}{B49}{P1-T4}{Solvable Lie algebra: exact global formula}{\proved{} (Lie's theorem \cf)}{thm:lie}
If the terms generate a solvable Lie algebra $\mathfrak g$ of dimension $d$, then for every $A\in\mathfrak g$ there are entire functions $g_1,\dots,g_d$, obtained by $d$ quadratures, with $e^{tA}=\prod_{i=1}^de^{g_i(t)X_i}$; hence $\Eint^{\mathrm{mult}}$ has a closed form with $d$ scalar functions.
\end{dimcard}
\begin{proof}
\step{1}{A flag of ideals}
By Lie's theorem there is a basis with $\mathfrak i_k=\spn\{X_k,\dots,X_d\}$ ideals \cf; $\mathrm{ad}_{X_j}$ acts on $\mathfrak i_k/\mathfrak i_{k+1}$ by a scalar $\lambda_{jk}$.
\step{2}{A triangular Wei--Norman system}
By B45(L5), $\mathrm{Ad}(P_{i-1})X_i=\sum_{l\ge i}M_{il}X_l$ with entire $M_{il}$ and $M_{ii}=\exp(\sum_{j<i}\lambda_{ji}g_j)\ne0$. The system $\sum_{i\le l}g_i'M_{il}=c_l$ is triangular: $g_1=c_1t$, and each $g_l$ is a quadrature of an entire function.
\whyworks{a solvable algebra can be climbed one ideal at a time, each step a scalar ODE.}
\end{proof}
\uses{B45.}

\begin{dimcard}{B50}{P1-P5}{Non-solvable: local only}{\proved}{}
\phantomsection\label{dim:P1-C1}%
For every finite-dimensional $\mathfrak g$ the Wei--Norman factorisation exists with analytic germs $g_i$ near $t=0$; globality can fail: in $\mathfrak{sl}_2$ with $A=E-F$ in Gauss order, $g_{21}(t)=-\tan t$. \textbf{(P1-C1)} The earlier claim ``infinite Lie algebra $\Rightarrow$ no finite factorisation of $\Eint^{\mathrm{mult}}$'' is false as stated \falsest: $\Eint^{\mathrm{mult}}=e^{-tT_2}e^{-tT_1}e^{t(T_1+T_2)}$ is always a product of three exponentials.
\end{dimcard}
\begin{proof}
\step{1}{Local existence}
The map $g'\mapsto\sum_ig_i'\mathrm{Ad}(P_{i-1}(g))X_i$ is given by an entire matrix $\Xi(g)$ with $\Xi(0)=I$; the analytic ODE $g'=\Xi(g)^{-1}c$ has a local solution.
\step{2}{The example}
Direct computation in $\mathrm{Mat}_2$: the Gauss coordinates of $e^{t(E-F)}$, a rotation, involve $\tan t$.
\step{3}{(C1)}
Definition of $\Eint^{\mathrm{mult}}$; the right question is A24.
\whyworks{coordinates of the second kind exist near the identity but can run into singularities.}
\end{proof}
\uses{B45.}

%=====================================================================
\section{Boundaries}

\begin{dimcard}{B51}{P1-P6}{A boundary creates an interaction of order one}{\proved}{prop:boundary}
On $L^2(0,L)$ with Dirichlet diffusion $A_1=D\partial^2$, inflow transport $A_2=-c\partial$ and the advection--diffusion generator $A$: $S_2(t)=0$ for $t\ge L/c$, while $S(t)$ is injective. So for $t\ge L/c$, $\Eint^{\mathrm{mult}}(t)=S(t)$: the relative splitting error is $1$ for every nonzero datum.
\end{dimcard}
\begin{proof}
\step{1}{The transport empties the interval}
$S_2(t)f(x)=f(x-ct)\mathbf 1_{x>ct}$, which is $0$ for $ct\ge L$.
\step{2}{The full flow is injective}
$A$ is similar, by $e^{-cx/2D}$, to a self-adjoint operator, so $S(t)$ is an analytic semigroup of self-adjoint type, injective.
\whyworks{the boundary condition belongs to the sum of the terms and cannot be shared between them.}
\end{proof}
\uses{none.}

\begin{dimcard}{B52}{OPA-T2}{Boundary interaction: sharp rate}{\proved}{}
\phantomsection\label{dim:P3-T3}\phantomsection\label{dim:OPF5A-T5}%
For advection--diffusion with Dirichlet conditions on $(0,L)$ and a datum $u_0\in C^{k+1}$ vanishing at the ends, with contact order $k$ (first nonzero derivative at an endpoint):
\[
\|E(t)\|=K_k\,t^{k/2+3/4}\big(1+O(t^{1/2})\big),\qquad
K_k=|c|\,N_{k-1}\,D^{k/2-1/4}\,\frac{\sqrt{u_0^{(k)}(0)^2+u_0^{(k)}(L)^2}}{(k-1)!},
\]
for the first-order interaction (OPA-T2) and, with $\tilde N_k=kN_{k-1}$, for the full $\Eint^{\mathrm{mult}}$ (OPF5A-T5). Uniformly: $\|E(t)\|\le1.301\,|c|D^{1/4}\|u_0'\|_\infty t^{5/4}$. Flat data give $O(t^p)$ for every $p$. \textbf{(P3-T3)} The first-order matrix elements are $V_{mn}=-4c\,mn/(L(m^2-n^2))$ for $m+n$ odd and $0$ otherwise (a reflection symmetry).
\end{dimcard}
\begin{proof}
\step{1}{A commutator identity}
The first-order interaction is $\int_0^tS_1(t-s)[\,\cdot\,]S_1(s)\dd s$ applied to $Vu_0$; the commutator of $V=-c\partial$ with the Dirichlet heat flow is a boundary term involving $u_0'$ at the ends \led{OPA-L2}.
\step{2}{Localisation at the ends}
Near each end the Dirichlet heat kernel is that of the half-line, and $u_0(y)=\gamma y^k/k!+O(y^{k+1})$.
\step{3}{Scaling}
With $x=\sqrt{Dt}\,w$ the boundary term becomes $t^{k/2+3/4}$ times a universal profile, whose $L^2$ norm is $N_{k-1}$; the corrections are $O(t^{1/2})$ relative.
\step{4}{Full interaction}
A conjugation reduces $\Eint^{\mathrm{mult}}$ to the same boundary terms plus a remainder $O(t^{k+5/4})$ and an exponentially small image term \led{OPF5A-L5}.
\whyworks{the boundary interaction is born where the datum meets the wall, and the heat flow sees the wall at distance $\sqrt{Dt}$.}
\end{proof}
\uses{A23; research files \texttt{work/OP-ANA} \S3, \texttt{work/OP-F5A} (L5, T5).}

%=====================================================================
\section{Where the interaction is continuous}

\begin{dimcard}{B53}{OPA-T3}{The exact frontier of discontinuity}{\proved}{}
\phantomsection\label{dim:OPA-P3}%
For an evolution $\tau-\sum_na_n(\xi)$ on a bounded rectangle: (b) if the terms are at least two distinct non-constant monomials, $u_0\mapsto\Eint^{\mathrm{add}}$ is unbounded from $(\Bclass\cap\ker Q,H^m)$ to $L^2$, for every $m$ and every base time. \textbf{(OPA-P3)} If the terms are affinely proportional, $a_n=\theta_na+\kappa_n$ with real $\theta_n$, and $t_0+\theta_n(t-t_0)\in\bar I$ (TC), then
\[
\Eint^{\mathrm{add}}(u)(t)=u(t)+(N-1)u(t_0)-\sum_ne^{(t-t_0)\kappa_n}u\big(t_0+\theta_n(t-t_0)\big),
\]
which is continuous in $H^1$; without (TC) it is unbounded. The general non-proportional case (e) is proved modulo the Tur\'an--Nazarov inequality \cf.
\end{dimcard}
\begin{proof}
\step{1}{Exponential test data}
On $u=e^{a(\zeta)t+\zeta x}$ with real $\zeta$ the interaction is $e_{\ge2}(z)\,u$, $z_n=e^{sa_n(\zeta)}-1$ (A19). If some term grows faster than the full symbol along a ray of $\zeta$ (condition (S)), the ratio to any $H^m$ norm grows exponentially, as in A21.
\step{2}{Proportional terms}
$e^{sa_n(\partial)}=e^{s\kappa_n}e^{s\theta_na(\partial)}$, and $e^{\tau a(\partial)}u_0=u(t_0+\tau)$ for $t_0+\tau\in\bar I$: the flows of the terms are time-rescaled copies of the true solution, controlled by $\|u\|_{H^1}$.
\whyworks{the interaction is continuous only when every term is a rescaled copy of the whole equation.}
\end{proof}
\uses{A19, A21, B43.}

\begin{dimcard}{B54}{P3-T6}{Balance decomposition}{\pw}{}
For $u\in H^m(\Omega)$ with $(\partial_t-\sum_ia_i(\partial_x))u=0$, $m\ge\max(\deg a_i,1)$, put $u^{(i)}=\int_{t_0}^ta_i(\partial_x)u(s)\dd s$. Then
\[
u=u_0\otimes1+\sum_iu^{(i)},\qquad \|u^{(i)}\|_{H^{m-d_i}}\le C\|u\|_{H^m}.
\]
In $\Bclass$ the $u^{(i)}$ solve a coupled Faddeev system whose Neumann series converges on every bounded time interval, with $\|\mathbb K^n\|\le(\kappa_0|I|)^n/n!$.
\end{dimcard}
\begin{proof}
\step{1}{Balance}
Integrate the equation in time: $u(t)-u_0=\sum_i\int_{t_0}^ta_i(\partial_x)u$. Each term loses $d_i$ derivatives in $x$ and gains one in $t$.
\step{2}{Faddeev system}
$\partial_tu^{(i)}=a_i(\partial)u=a_i(\partial)(u_0+\sum_ju^{(j)})$, so $u^{(i)}=G_ia_i(\partial)(u_0+\sum_{j\ne i}u^{(j)})$ with the Green operator $G_i$ of $\partial_t-a_i$; the off-diagonal Volterra operators have the factorial bound.
\whyworks{a split that uses the true solution, not the single-term flows, inherits its continuity.}
\end{proof}
\uses{B44.}

%=====================================================================
\section{Regularisation of the term split}

\begin{dimcard}{B55}{P2-T1}{Every regulariser has the same limit residual}{\proved}{}
\phantomsection\label{dim:P2-L1}%
Let $S:W\to H$ be linear with closure of the range $M$, $R$ any regulariser with $S(\mathrm{dom}\,R)$ dense in $M$, and $\omega_\alpha$ near-minimisers of $\tfrac12\|S\omega-u\|^2+\alpha R(\omega)$ (errors $\varepsilon_\alpha\to0$). Then $u-S\omega_\alpha\to P_M^\perp u$, so the limit residual is $\mathrm{dist}(u,M)$ for every regulariser. \textbf{(P2-L1)} If $Q$ kills $M$, then $\|E^\perp(u)\|_{H^m}\ge\|Q(\partial)u\|_{H^{m-k}}/C_Q$.
\end{dimcard}
\begin{proof}
\step{1}{Upper bound}
For any $v\in\mathrm{dom}R$, $\tfrac12\|S\omega_\alpha-u\|^2\le\tfrac12\|Sv-u\|^2+\alpha R(v)+\varepsilon_\alpha$; let $\alpha\to0$, then choose $Sv$ close to $P_Mu$.
\step{2}{Lower bound}
$\|u-S\omega\|\ge\mathrm{dist}(u,M)$ always; Pythagoras in $u=P_Mu+P_M^\perp u$ gives norm convergence.
\step{3}{(L1)}
$Q(\partial)u=Q(\partial)E^\perp(u)$ and $Q(\partial)$ is bounded $H^m\to H^{m-k}$.
\whyworks{a regulariser only chooses among representations; it cannot reach what lies outside their closure.}
\end{proof}
\uses{none.}

\begin{dimcard}{B56}{P2-P3}{The regulariser acts only on the ambiguity}{\proved}{}
\phantomsection\label{dim:P2-P2}%
The exact representations of $u$ form $\omega^\dagger+\ker S$. For convex, lower semicontinuous, coercive $R$ the regularised solutions converge (along subsequences) to an $R$-minimal exact representation, to $S^\dagger u$ for $R=\tfrac12\|\cdot\|^2$. On a rectangle $\ker S$ has dimension $\delta(P)$ for factors and $\sum_{n<m}\deg(a_n-a_m)$ for terms. \textbf{(P2-P2)} For ODEs, the sum of the per-term kernels is $\ker\mathrm{lcm}(p_n)(D)$.
\end{dimcard}
\begin{proof}
\step{1}{Bounded minimisers}
$R(\omega_\alpha)\le R(\omega^\dagger)+\varepsilon_\alpha/\alpha$ and coercivity bound $\omega_\alpha$; weak limits are exact and $R$-minimal by weak lower semicontinuity (Mazur).
\step{2}{Dimensions}
A13(ii) for factors; for terms $\ring/(\tau-a_n,\tau-a_m)\cong\C[\xi]/(a_n-a_m)$.
\step{3}{(P2)}
The roots of the lcm, with maximal multiplicities, give a basis contained in the union of the per-term kernels.
\whyworks{all methods agree up to a finite-dimensional gauge, where the regulariser picks a representative.}
\end{proof}
\uses{A13.}

\begin{dimcard}{B57}{P2-P4}{Linear encoders do not help}{\proved}{}
(a) A bounded, boundedly invertible linear encoder $E$ preserves closedness of sums and distances up to the constants $m_E$, $\|E\|$. (b) An injective compact encoder can make a non-dense $M$ have dense image, but then the encoded components diverge.
\end{dimcard}
\begin{proof}
\step{1}{(a)}
$m_E\|u-v\|\le\|Eu-Ev\|\le\|E\|\|u-v\|$; isomorphisms preserve closed subspaces.
\step{2}{(b)}
With $Ee_k=e_k/k$ and $M=w^\perp$, $w\propto\sum_ke_k/k$: $y\perp EM$ forces $Ey\in\C w$, impossible in $\ell^2$ unless $y=0$. If $Em_k\to Eu$ with $u\notin M$ and $m_k$ bounded, a weak limit would give $u\in M$.
\whyworks{stable encodings preserve geometry; unstable ones hide it by sending pieces to infinity.}
\end{proof}
\uses{none.}

\begin{dimcard}{B58}{P2-P5}{Where the interaction is stable}{\proved}{}
(a) In $\Bclass_\rho$: $\sup_t\|\Eint(t)\|_{B_{\rho'}}\le e_{\ge2}(\beta)\|u_0\|_{B_{\rho'}}$ with $\beta_n=e^{|I||a_n|(\rho')}-1$. (b) On $\R$ with $L^2$ data: $\|\Eint(t)\|_{L^2}\le e_{\ge2}(\beta)\|u_0\|$ with $\beta_n=\sup|e^{ta_n(i\kappa)}-1|$; for advection--diffusion $\|\Eint\|\le2\|u_0\|$.
\end{dimcard}
\begin{proof}
\step{1}{(a)}
A19(ii) and $\|e^{sa_n}-I\|\le e^{|s||a_n|(\rho')}-1$; $e_{\ge2}$ has nonnegative coefficients.
\step{2}{(b)}
Plancherel: $\widehat{\Eint}=e_{\ge2}(z)\hat u_0$ and $|e_{\ge2}(z)|\le e_{\ge2}(|z|)$. For advection--diffusion $|z_1|\le2$, $|z_2|\le1$.
\whyworks{stability of the interaction is a property of the data class, not of the method.}
\end{proof}
\uses{A19, B44.}

\begin{dimcard}{B59}{P3-C1}{The hypotheses of Dyson--Phillips cannot be dropped}{\proved}{}
(a) Same order ($A_0=\partial^2$, $V=-2\partial^2$): the Dyson terms have norm $(2n)^ne^{-n}/n!\sim2^n/\sqrt{2\pi n}$ and $A_0+V$ generates no semigroup. (b) Dispersive principal term ($A_0=i\partial^2$, $V=i\partial$): $A_0+V$ generates no semigroup. (c) $A_0=i\partial^2$, $V=\partial$: the sum generates a group but the Dyson terms are unbounded.
\end{dimcard}
\begin{proof}
\step{1}{Multipliers}
All operators are Fourier multipliers; in (a) $\sup_y(2ty)^ne^{-ty}/n!=(2n)^ne^{-n}/n!$.
\step{2}{Non-generation}
The spectrum of a multiplier is the closure of its symbol's range; in (a), (b) the real part is unbounded above, so by Hille--Yosida there is no generator.
\whyworks{without smoothing by the principal term, each perturbation step costs a full derivative.}
\end{proof}
\uses{A23.}

\begin{dimcard}{B60}{OPF-T2}{Repeated factors: the multiplicity term}{\pw}{}
\phantomsection\label{dim:OPF-T1}%
For $Q=(\tau-a(\xi))^e$ and $u\in\Bclass$, $U=(u^{(0)},\dots,u^{(e-1)})$ with $u^{(k)}=(\partial_t-a)^ku$ satisfies $U(t)=e^{s(a(\partial_x)I+N_e)}F$ with the nilpotent shift $N_e$; the multiplicity term commutes with the others, its flow is polynomial in $s$ (secular growth), and $\Eint^{\mathrm{mult}}=0$. \textbf{(OPF-T1)} For every $P$, every factor decomposition of $u\in\Bclass$ lies in $\Bclass$.
\end{dimcard}
\begin{proof}
\step{1}{(T2)}
$\partial_tU=(a(\partial_x)I+N_e)U$, a system with commuting constant-coefficient generator; B45(L3) for systems.
\step{2}{(T1)}
For each $j$, B26(c) gives $r(\partial_x)u_j$ and $s(\partial_t)u_j$ in $\Bclass$; B45(L4) inverts $r(\partial_x)$ inside $\Bclass$, and the remaining finite-dimensional part is controlled as in B47.
\whyworks{a repeated factor adds a nilpotent term, which commutes with everything and grows polynomially.}
\end{proof}
\uses{B26, B45, B47.}

\section{Results used without proof here}

\listcard{TROP-3$'$, TROP-4}{Uniform-in-time bound of the interaction on dissipative faces of the Newton polygon; no uniform bound on conservative faces.}{\proved}{Proofs in \texttt{work/SCOUT-TROP/VERDETTO.md} (audited); not reproduced. TROP-5..8: \pw.}
\listcard{R5-P6, R5-T4}{Global Wei--Norman for Ornstein--Uhlenbeck; Duhamel with bound.}{\proved}{Proofs in \texttt{work/CHECK/DECLASSAMENTI.md} \S B; not reproduced.}
\listcard{P3-T2, P3-T4, P3-T5}{Path formula for Dyson terms; closed form for an $\varepsilon x$ perturbation; Appell and Rosenbloom--Widder duals.}{\proved{} / \pw}{Proofs in \texttt{work/P3}; not reproduced.}
\listcard{OPA-P4}{Splitting error at full coupling equals the first-order term to relative $O(t^{1/2})$.}{\proved}{Now contained in B52 (OPF5A-T5).}
\listcard{P1-Q6}{No-go of continuity for every operator with at least two terms.}{closed}{By B53 (monomials) and B53(e) mod Tur\'an--Nazarov.}
\listcard{P1-Q3, P1-Q7}{No finite product of flows for $\partial^2$ and $x^k$, $k\ge3$; A24 without the switching-off hypothesis.}{\openq}{}
